% Options for packages loaded elsewhere
% Options for packages loaded elsewhere
\PassOptionsToPackage{unicode}{hyperref}
\PassOptionsToPackage{hyphens}{url}
%
\documentclass[
  ignorenonframetext,
  aspectratio=169,
]{beamer}
\newif\ifbibliography
\usepackage{pgfpages}
\setbeamertemplate{caption}[numbered]
\setbeamertemplate{caption label separator}{: }
\setbeamercolor{caption name}{fg=normal text.fg}
\beamertemplatenavigationsymbolsframe
% remove section numbering
\setbeamertemplate{part page}{
  \centering
  \begin{beamercolorbox}[sep=16pt,center]{part title}
    \usebeamerfont{part title}\insertpart\par
  \end{beamercolorbox}
}
\setbeamertemplate{section page}{
  \centering
  \begin{beamercolorbox}[sep=12pt,center]{section title}
    \usebeamerfont{section title}\insertsection\par
  \end{beamercolorbox}
}
\setbeamertemplate{subsection page}{
  \centering
  \begin{beamercolorbox}[sep=8pt,center]{subsection title}
    \usebeamerfont{subsection title}\insertsubsection\par
  \end{beamercolorbox}
}
% Prevent slide breaks in the middle of a paragraph
\widowpenalties 1 10000
\raggedbottom
\AtBeginPart{
  \frame{\partpage}
}
\AtBeginSection{
  \ifbibliography
  \else
    \frame{\sectionpage}
  \fi
}
\AtBeginSubsection{
  \frame{\subsectionpage}
}
\usepackage{iftex}
\ifPDFTeX
  \usepackage[T1]{fontenc}
  \usepackage[utf8]{inputenc}
  \usepackage{textcomp} % provide euro and other symbols
\else % if luatex or xetex
  \usepackage{unicode-math} % this also loads fontspec
  \defaultfontfeatures{Scale=MatchLowercase}
  \defaultfontfeatures[\rmfamily]{Ligatures=TeX,Scale=1}
\fi
\usepackage{lmodern}

\usetheme[]{Madrid}
\usecolortheme[]{default}
\usefonttheme[]{professionalfonts}
\ifPDFTeX\else
  % xetex/luatex font selection
\fi
% Use upquote if available, for straight quotes in verbatim environments
\IfFileExists{upquote.sty}{\usepackage{upquote}}{}
\IfFileExists{microtype.sty}{% use microtype if available
  \usepackage[]{microtype}
  \UseMicrotypeSet[protrusion]{basicmath} % disable protrusion for tt fonts
}{}
\makeatletter
\@ifundefined{KOMAClassName}{% if non-KOMA class
  \IfFileExists{parskip.sty}{%
    \usepackage{parskip}
  }{% else
    \setlength{\parindent}{0pt}
    \setlength{\parskip}{6pt plus 2pt minus 1pt}}
}{% if KOMA class
  \KOMAoptions{parskip=half}}
\makeatother

\usepackage{color}
\usepackage{fancyvrb}
\newcommand{\VerbBar}{|}
\newcommand{\VERB}{\Verb[commandchars=\\\{\}]}
\DefineVerbatimEnvironment{Highlighting}{Verbatim}{commandchars=\\\{\}}
% Add ',fontsize=\small' for more characters per line
\usepackage{framed}
\definecolor{shadecolor}{RGB}{241,243,245}
\newenvironment{Shaded}{\begin{snugshade}}{\end{snugshade}}
\newcommand{\AlertTok}[1]{\textcolor[rgb]{0.68,0.00,0.00}{#1}}
\newcommand{\AnnotationTok}[1]{\textcolor[rgb]{0.37,0.37,0.37}{#1}}
\newcommand{\AttributeTok}[1]{\textcolor[rgb]{0.40,0.45,0.13}{#1}}
\newcommand{\BaseNTok}[1]{\textcolor[rgb]{0.68,0.00,0.00}{#1}}
\newcommand{\BuiltInTok}[1]{\textcolor[rgb]{0.00,0.23,0.31}{#1}}
\newcommand{\CharTok}[1]{\textcolor[rgb]{0.13,0.47,0.30}{#1}}
\newcommand{\CommentTok}[1]{\textcolor[rgb]{0.37,0.37,0.37}{#1}}
\newcommand{\CommentVarTok}[1]{\textcolor[rgb]{0.37,0.37,0.37}{\textit{#1}}}
\newcommand{\ConstantTok}[1]{\textcolor[rgb]{0.56,0.35,0.01}{#1}}
\newcommand{\ControlFlowTok}[1]{\textcolor[rgb]{0.00,0.23,0.31}{\textbf{#1}}}
\newcommand{\DataTypeTok}[1]{\textcolor[rgb]{0.68,0.00,0.00}{#1}}
\newcommand{\DecValTok}[1]{\textcolor[rgb]{0.68,0.00,0.00}{#1}}
\newcommand{\DocumentationTok}[1]{\textcolor[rgb]{0.37,0.37,0.37}{\textit{#1}}}
\newcommand{\ErrorTok}[1]{\textcolor[rgb]{0.68,0.00,0.00}{#1}}
\newcommand{\ExtensionTok}[1]{\textcolor[rgb]{0.00,0.23,0.31}{#1}}
\newcommand{\FloatTok}[1]{\textcolor[rgb]{0.68,0.00,0.00}{#1}}
\newcommand{\FunctionTok}[1]{\textcolor[rgb]{0.28,0.35,0.67}{#1}}
\newcommand{\ImportTok}[1]{\textcolor[rgb]{0.00,0.46,0.62}{#1}}
\newcommand{\InformationTok}[1]{\textcolor[rgb]{0.37,0.37,0.37}{#1}}
\newcommand{\KeywordTok}[1]{\textcolor[rgb]{0.00,0.23,0.31}{\textbf{#1}}}
\newcommand{\NormalTok}[1]{\textcolor[rgb]{0.00,0.23,0.31}{#1}}
\newcommand{\OperatorTok}[1]{\textcolor[rgb]{0.37,0.37,0.37}{#1}}
\newcommand{\OtherTok}[1]{\textcolor[rgb]{0.00,0.23,0.31}{#1}}
\newcommand{\PreprocessorTok}[1]{\textcolor[rgb]{0.68,0.00,0.00}{#1}}
\newcommand{\RegionMarkerTok}[1]{\textcolor[rgb]{0.00,0.23,0.31}{#1}}
\newcommand{\SpecialCharTok}[1]{\textcolor[rgb]{0.37,0.37,0.37}{#1}}
\newcommand{\SpecialStringTok}[1]{\textcolor[rgb]{0.13,0.47,0.30}{#1}}
\newcommand{\StringTok}[1]{\textcolor[rgb]{0.13,0.47,0.30}{#1}}
\newcommand{\VariableTok}[1]{\textcolor[rgb]{0.07,0.07,0.07}{#1}}
\newcommand{\VerbatimStringTok}[1]{\textcolor[rgb]{0.13,0.47,0.30}{#1}}
\newcommand{\WarningTok}[1]{\textcolor[rgb]{0.37,0.37,0.37}{\textit{#1}}}

\usepackage{longtable,booktabs,array}
\newcounter{none} % for unnumbered tables
\usepackage{calc} % for calculating minipage widths
\usepackage{caption}
% Make caption package work with longtable
\makeatletter
\def\fnum@table{\tablename~\thetable}
\makeatother
\usepackage{graphicx}
\makeatletter
\newsavebox\pandoc@box
\newcommand*\pandocbounded[1]{% scales image to fit in text height/width
  \sbox\pandoc@box{#1}%
  \Gscale@div\@tempa{\textheight}{\dimexpr\ht\pandoc@box+\dp\pandoc@box\relax}%
  \Gscale@div\@tempb{\linewidth}{\wd\pandoc@box}%
  \ifdim\@tempb\p@<\@tempa\p@\let\@tempa\@tempb\fi% select the smaller of both
  \ifdim\@tempa\p@<\p@\scalebox{\@tempa}{\usebox\pandoc@box}%
  \else\usebox{\pandoc@box}%
  \fi%
}
% Set default figure placement to htbp
\def\fps@figure{htbp}
\makeatother





\setlength{\emergencystretch}{3em} % prevent overfull lines

\providecommand{\tightlist}{%
  \setlength{\itemsep}{0pt}\setlength{\parskip}{0pt}}



 


\makeatletter
\@ifpackageloaded{caption}{}{\usepackage{caption}}
\AtBeginDocument{%
\ifdefined\contentsname
  \renewcommand*\contentsname{Table of contents}
\else
  \newcommand\contentsname{Table of contents}
\fi
\ifdefined\listfigurename
  \renewcommand*\listfigurename{List of Figures}
\else
  \newcommand\listfigurename{List of Figures}
\fi
\ifdefined\listtablename
  \renewcommand*\listtablename{List of Tables}
\else
  \newcommand\listtablename{List of Tables}
\fi
\ifdefined\figurename
  \renewcommand*\figurename{Figure}
\else
  \newcommand\figurename{Figure}
\fi
\ifdefined\tablename
  \renewcommand*\tablename{Table}
\else
  \newcommand\tablename{Table}
\fi
}
\@ifpackageloaded{float}{}{\usepackage{float}}
\floatstyle{ruled}
\@ifundefined{c@chapter}{\newfloat{codelisting}{h}{lop}}{\newfloat{codelisting}{h}{lop}[chapter]}
\floatname{codelisting}{Listing}
\newcommand*\listoflistings{\listof{codelisting}{List of Listings}}
\makeatother
\makeatletter
\makeatother
\makeatletter
\@ifpackageloaded{caption}{}{\usepackage{caption}}
\@ifpackageloaded{subcaption}{}{\usepackage{subcaption}}
\makeatother

\usepackage{bookmark}
\IfFileExists{xurl.sty}{\usepackage{xurl}}{} % add URL line breaks if available
\urlstyle{same}
\hypersetup{
  pdftitle={Sesión 9: Pruebas Diagnósticas I},
  pdfauthor={Jesús Esteban Hernández},
  hidelinks,
  pdfcreator={LaTeX via pandoc}}


\title{\texorpdfstring{Sesión 9: Pruebas Diagnósticas
I}{Sesión 9: Pruebas Diagnósticas I}}
\subtitle{\texorpdfstring{Teorema de Bayes: validación y rendimiento de
una prueba
diagnóstica}{Teorema de Bayes: validación y rendimiento de una prueba diagnóstica}}
\author{Jesús Esteban Hernández}
\date{2026-10-12}

\begin{document}
\frame{\titlepage}


\begin{frame}
\#\# Objetivos de aprendizaje \{.allowframebreaks\}

\begin{itemize}
\tightlist
\item
  Comprender los fundamentos de la \textbf{validación diagnóstica}.
\item
  Construir e interpretar \textbf{tablas de contingencia
  \(2 \times 2\)}.
\item
  Calcular e interpretar \textbf{sensibilidad, especificidad, VPP y
  VPN}.
\item
  Analizar el impacto de la \textbf{prevalencia} mediante el
  \textbf{Teorema de Bayes}.
\end{itemize}
\end{frame}

\begin{frame}
\# Repaso de probabilidad

\#\# Experimento Aleatorio y Espacio Muestral

\textbf{Experimento Aleatorio}:

\begin{itemize}
\tightlist
\item
  Repetible en las mismas condiciones.
\item
  Resultados distintos con mismas condiciones.
\item
  Todos los resultados conocidos \emph{a priori}.
\item
  No se puede predecir el resultado particular.
\end{itemize}

\textbf{Espacio Muestral (\(\Omega\))}: todos los resultados posibles.
\end{frame}

\begin{frame}
\#\# Tipos de Espacios Muestrales

\begin{tabular}{llll}
\toprule
Tipo & Ejemplo & $\Omega$ & Suceso \\
\midrule
Finito & Grupo sanguíneo & $\{A,B,AB,O\}$ & Compatibilidad \\
Numerable & Crisis/año & $\{0,1,2,\dots\}$ & $>3$ crisis \\
Continuo & Glucosa & $(0,+\infty)$ & Diabetes \\
\bottomrule
\end{tabular}
\end{frame}

\begin{frame}
\#\# Espacio Probabilístico \((\Omega, \mathcal{A}, P)\)

\begin{enumerate}
\tightlist
\item
  \(\Omega\): espacio muestral.
\item
  \(\mathcal{A}\): \(\sigma\)-álgebra de sucesos.
\item
  \(P\): medida de probabilidad \(P: \mathcal{A} \to [0,1]\).
\end{enumerate}
\end{frame}

\begin{frame}
\#\# Axiomas de Kolmogorov

\begin{itemize}
\tightlist
\item
  \textbf{A1}: \(P(A) \ge 0\)
\item
  \textbf{A2}: \(P(\Omega) = 1\)
\item
  \textbf{A3}: Si \(A_i \cap A_j = \emptyset\):
  \[P\left(\bigcup_{i=1}^{\infty} A_i\right) = \sum_{i=1}^{\infty} P(A_i)\]
\end{itemize}
\end{frame}

\begin{frame}
\#\# Propiedades Fundamentales

\begin{enumerate}
\tightlist
\item
  \(P(A') = 1 - P(A)\)
\item
  \(P(\emptyset) = 0\)
\item
  \(A \subset B \implies P(A) \le P(B)\)
\item
  \(P(A \cup B) = P(A) + P(B) - P(A \cap B)\)
\end{enumerate}
\end{frame}

\begin{frame}
\end{frame}

\subsection{Independencia}\label{independencia}

\begin{frame}{Independencia}
\(A \perp B \iff P(A \cap B) = P(A)P(B)\)

\textbf{Ojo}: disjuntos \(\ne\) independientes.
\end{frame}

\begin{frame}
\end{frame}

\subsection{Probabilidad Condicionada}\label{probabilidad-condicionada}

\begin{frame}{Probabilidad Condicionada}
\[P(B \mid A) = \frac{P(A \cap B)}{P(A)}, \quad P(A) > 0\]

\textbf{Ley multiplicativa}:

\[P(A \cap B) = P(A) \cdot P(B \mid A)\]
\end{frame}

\begin{frame}
\end{frame}

\subsection{Sistema Completo de
Sucesos}\label{sistema-completo-de-sucesos}

\begin{frame}{Sistema Completo de Sucesos}
Partición de \(\Omega\): sucesos incompatibles dos a dos cuya unión es
\(\Omega\).
\end{frame}

\begin{frame}
\end{frame}

\section{Teorema de Bayes}\label{teorema-de-bayes}

\subsection{Probabilidad Total}\label{probabilidad-total}

\begin{frame}{Probabilidad Total}
\[P(B) = \sum_{i=1}^{k} P(A_i) P(B \mid A_i)\]
\end{frame}

\subsection{Teorema de Bayes}\label{teorema-de-bayes-1}

\begin{frame}{Teorema de Bayes}
\[P(A_j \mid B) = \frac{P(A_j) P(B \mid A_j)}{\sum_i P(A_i) P(B \mid A_i)}\]
\end{frame}

\begin{frame}
\end{frame}

\subsection{Bayes aplicado al
diagnóstico}\label{bayes-aplicado-al-diagnuxf3stico}

\begin{frame}{Bayes aplicado al diagnóstico}
\[P(E \mid S) = \frac{P(E) P(S \mid E)}{P(E)P(S\mid E) + P(\bar{E})P(S \mid \bar{E})}\]

\begin{enumerate}
\tightlist
\item
  Depende de la previa \(P(E)\).
\item
  El síntoma debe \textbf{no ser independiente} de la enfermedad.
\item
  Los síntomas útiles \textbf{discriminan}, no son los más frecuentes.
\end{enumerate}
\end{frame}

\begin{frame}
\end{frame}

\section{Validación Diagnóstica}\label{validaciuxf3n-diagnuxf3stica}

\subsection{\texorpdfstring{Matriz de Confusión
\(2 \times 2\)}{Matriz de Confusión 2 \textbackslash times 2}}\label{matriz-de-confusiuxf3n-2-times-2}

\begin{frame}{Matriz de Confusión \(2 \times 2\)}
\begin{tabular}{lccc}
\toprule
 & $D+$ & $D-$ & Total \\
\midrule
$T+$ & $VP$ & $FP$ & $VP+FP$ \\
$T-$ & $FN$ & $VN$ & $FN+VN$ \\
\midrule
Total & $VP+FN$ & $FP+VN$ & $N$ \\
\bottomrule
\end{tabular}
\end{frame}

\begin{frame}
\end{frame}

\subsection{Sensibilidad y
Especificidad}\label{sensibilidad-y-especificidad}

\begin{frame}{Sensibilidad y Especificidad}
\[s = \frac{VP}{VP+FN} \qquad e = \frac{VN}{FP+VN}\]

\emph{No dependen de la prevalencia} (en teoría).
\end{frame}

\begin{frame}
\end{frame}

\subsection{Valores Predictivos}\label{valores-predictivos}

\begin{frame}{Valores Predictivos}
\[\text{VPP} = \frac{VP}{VP+FP} \qquad \text{VPN} = \frac{VN}{FN+VN}\]

\emph{Sí dependen de la prevalencia.}
\end{frame}

\begin{frame}
\end{frame}

\subsection{Factor de Bayes}\label{factor-de-bayes}

\begin{frame}{Factor de Bayes}
\[Odd(E \mid S) = Odd(E) \cdot \frac{P(S \mid E)}{P(S \mid \bar{E})} = Odd(E) \cdot \text{FB}\]

En función de \(s\) y \(e\):

\[Odd(E \mid S) = Odd(E) \cdot \frac{s}{1-e}\]
\end{frame}

\begin{frame}
\end{frame}

\subsection{Impacto de la Prevalencia}\label{impacto-de-la-prevalencia}

\begin{frame}{Impacto de la Prevalencia}
\[\text{VPP} = \frac{s \cdot p}{s \cdot p + (1-e)(1-p)}\]

\textbf{En cribado} (prevalencia baja) el VPP cae drásticamente.
\end{frame}

\begin{frame}
\end{frame}

\section{Caso práctico en R}\label{caso-pruxe1ctico-en-r}

\subsection{Datos}\label{datos}

\begin{frame}[fragile]{Datos}
\texttt{\_data/preven\_cc.csv}:

\begin{itemize}
\tightlist
\item
  \texttt{gold\_std}: Sano / Enfermo
\item
  \texttt{biomarcador}: continuo
\end{itemize}

Punto de corte: \textbf{52}.
\end{frame}

\begin{frame}
\end{frame}

\subsection{Código}\label{cuxf3digo}

\begin{frame}[fragile]{Código}
\begin{Shaded}
\begin{Highlighting}[]
\FunctionTok{library}\NormalTok{(epiR)}
\NormalTok{biomark1 }\OtherTok{\textless{}{-}} \FunctionTok{read.csv}\NormalTok{(}\StringTok{"\_data/preven\_cc.csv"}\NormalTok{)}
\NormalTok{biomark1}\SpecialCharTok{$}\NormalTok{resultado\_prueba }\OtherTok{\textless{}{-}} \FunctionTok{factor}\NormalTok{(}
  \FunctionTok{ifelse}\NormalTok{(biomark1}\SpecialCharTok{$}\NormalTok{biomarcador }\SpecialCharTok{\textgreater{}} \DecValTok{52}\NormalTok{, }\DecValTok{1}\NormalTok{, }\DecValTok{0}\NormalTok{),}
  \AttributeTok{levels =} \FunctionTok{c}\NormalTok{(}\DecValTok{0}\NormalTok{,}\DecValTok{1}\NormalTok{), }\AttributeTok{labels =} \FunctionTok{c}\NormalTok{(}\StringTok{"Negativo"}\NormalTok{,}\StringTok{"Positivo"}\NormalTok{)}
\NormalTok{)}

\NormalTok{VP }\OtherTok{\textless{}{-}} \FunctionTok{sum}\NormalTok{(biomark1}\SpecialCharTok{$}\NormalTok{resultado\_prueba}\SpecialCharTok{==}\StringTok{"Positivo"} \SpecialCharTok{\&}\NormalTok{ biomark1}\SpecialCharTok{$}\NormalTok{gold\_std}\SpecialCharTok{==}\StringTok{"Enfermo"}\NormalTok{)}
\NormalTok{FP }\OtherTok{\textless{}{-}} \FunctionTok{sum}\NormalTok{(biomark1}\SpecialCharTok{$}\NormalTok{resultado\_prueba}\SpecialCharTok{==}\StringTok{"Positivo"} \SpecialCharTok{\&}\NormalTok{ biomark1}\SpecialCharTok{$}\NormalTok{gold\_std}\SpecialCharTok{==}\StringTok{"Sano"}\NormalTok{)}
\NormalTok{FN }\OtherTok{\textless{}{-}} \FunctionTok{sum}\NormalTok{(biomark1}\SpecialCharTok{$}\NormalTok{resultado\_prueba}\SpecialCharTok{==}\StringTok{"Negativo"} \SpecialCharTok{\&}\NormalTok{ biomark1}\SpecialCharTok{$}\NormalTok{gold\_std}\SpecialCharTok{==}\StringTok{"Enfermo"}\NormalTok{)}
\NormalTok{VN }\OtherTok{\textless{}{-}} \FunctionTok{sum}\NormalTok{(biomark1}\SpecialCharTok{$}\NormalTok{resultado\_prueba}\SpecialCharTok{==}\StringTok{"Negativo"} \SpecialCharTok{\&}\NormalTok{ biomark1}\SpecialCharTok{$}\NormalTok{gold\_std}\SpecialCharTok{==}\StringTok{"Sano"}\NormalTok{)}
\end{Highlighting}
\end{Shaded}
\end{frame}

\begin{frame}
\end{frame}

\subsection{Matriz y s, e}\label{matriz-y-s-e}

\begin{frame}[fragile]{Matriz y s, e}
\begin{verbatim}
      Enfermo Sano
Test+     137  408
Test-      20  435
\end{verbatim}

\begin{verbatim}
 sensibilidad especificidad 
    0.8726115     0.5160142 
\end{verbatim}
\end{frame}

\begin{frame}
\end{frame}

\subsection{\texorpdfstring{\texttt{epi.tests}}{epi.tests}}\label{epi.tests}

\begin{frame}[fragile]{\texttt{epi.tests}}
\begin{verbatim}
   statistic         est       lower       upper
1         ap  0.45500000  0.42380745  0.48645821
2         tp  0.15700000  0.13498766  0.18105255
3         se  0.12738854  0.07957909  0.18986556
4         sp  0.48398577  0.44976414  0.51831970
5    diag.ac  0.42800000  0.39708390  0.45934120
6    diag.or  0.13692424  0.08403282  0.22310624
7       nndx -2.57317002 -3.42683175 -2.12469057
8     youden -0.38862570 -0.47065677 -0.29181474
9     pv.pos  0.04395604  0.02705301  0.06707227
10    pv.neg  0.74862385  0.70999733  0.78452993
11    lr.pos  0.24687020  0.16308632  0.37369715
12    lr.neg  1.80296928  1.64479815  1.97635085
13    p.rout  0.54500000  0.51354179  0.57619255
14     p.rin  0.45500000  0.42380745  0.48645821
15    p.tpdn  0.51601423  0.48168030  0.55023586
16    p.tndp  0.87261146  0.81013444  0.92042091
17    p.dntp  0.95604396  0.93292773  0.97294699
18    p.dptn  0.25137615  0.21547007  0.29000267
\end{verbatim}
\end{frame}

\begin{frame}
\end{frame}

\subsection{Discusión}\label{discusiuxf3n}

\begin{frame}{Discusión}
\begin{enumerate}
\tightlist
\item
  ¿Qué pasa con el VPP si prevalencia = 1\%?
\item
  ¿A quién aplicaríamos el test?
\item
  ¿Y si el corte fuese 45?
\end{enumerate}
\end{frame}

\begin{frame}
\end{frame}

\subsection{Referencias}\label{referencias}

\begin{frame}{Referencias}
\begin{itemize}
\tightlist
\item
  Pepe, M.S. (2003). \emph{The Statistical Evaluation of Medical Tests}.
  Oxford.
\item
  Zhou, X.H. et al.~(2011). \emph{Statistical Methods in Diagnostic
  Medicine}. Wiley.
\item
  Stevenson M. et al.~(2024). \emph{epiR}.
\end{itemize}
\end{frame}




\end{document}
